\begin{center}
{\heiti\zihao{4}中国研究生创新实践系列大赛\par}
\vspace{.4em}
{\heiti\zihao{4}“华为杯”第二十三届中国研究生\par 数学建模竞赛\par}
\vspace{1em}
{\heiti\zihao{3}题目：复杂场景下多模态情感预测的\par 时序建模与可解释融合\par}
\vspace{.7em}
{\heiti\zihao{4}摘要}
\end{center}
针对多模态情感预测中的时间对应、局部信息缺失及决策解释，建立有效观测约束下的特征组织、联合预测与可加解释方法。实验使用题目附件，三项任务分别输出全量特征、专项预测和局部证据。

问题一按实际音视频时钟提取上下文化文本、声学量、画面与人脸形变特征，分离物理时间支持、逐维有效性和内容对应。利用无提示语音识别与双上下文音素对齐生成候选，以词序和时间双重约束求解词组选择。100条样本全部保留；1932个原词中1397个获得词级候选，50个词组另覆盖71个词，总候选覆盖率75.98\%。固定证据下，动态规划保持覆盖、减少3个词组并消除2对不兼容关系。同环境重建的3100个特征与索引数组一致。

问题二采用编码前缺失屏蔽、观测掩码汇总和可用比例输入，联合优化极性分类与强度回归。通过视频分组三折选择逆频率类别权重，在验证集三种子确认中，Macro-F1均值由0.5959提高到0.6163，中性F1由0.3871提高到0.4808，Accuracy由0.6369降至0.6277，MAE为0.5790。连续缺失实验显示，文本缺失影响更大，缺失位置与片段形状也改变结果。给出附件三全部30条预测。

问题三用单模态项、成对交互和残差方差驱动的缩放构造融合模型，将分类间隔与回归预激活精确分解到模态和局部位置。验证集Accuracy为65.11\%、Macro-F1为0.6323、MAE为0.5397、Pearson为0.6993。验证、测试及专项样本共1475条的分类分解最大绝对误差为$1.52\times10^{-6}$。附件四20条均输出预测、贡献及关键位置，18条取得可回看的关键时间范围。

上述方法分别以候选覆盖与稳定性、预测误差、分解一致性评价时间组织、情感预测和模型解释，保留输入缺失及定位支持不足的明确状态。

\noindent\textbf{关键词：}多模态情感预测；多粒度时序对应；局部模态缺失；类别平衡；可加解释
